%% The following is a directive for TeXShop to indicate the main file
%%!TEX root = diss.tex

%=======================================
\section{Miscellaneous Things}
%=======================================

If we could prove the validity of one of the functional equations in the corollaries (5th row in each corollary) we would then be able to extract coefficients and obtain the formulas in the first corollary. However, since we have a direct proof already, this section is kind of irrelevant now. We'll just keep it for now because it has some interesting math. The steps from now on are almost exactly the same as those that you'd have to do to prove the Dyck path formula. We'll prove the last one involving a $b_1$ coefficient. First we have
\begin{align}
   K &= 1 - tS = 1 - t(x_1+\ol{x_1})(x_2+\ol{x_2})\\
   K_1 &= c_1 - t \ol{x_1} S_1^{-1} = c_1 + \wh{K}_1\\
   K_2 &= c_2 - t \ol{x_2} S_2^{-1}\\
   K_{1,2} &= c_1c_2 - t\ol{x_1x_2} = c_1c_2+\wh{K}_{1,2}
\end{align}
We also need the coefficient extractions
\begin{align}[x_1^0]:&\quad Q_1^0 - c_2Q_{1,2}^{0,0}-ta_1S_1^{-1}Q_1^1 =\ol{a_2}\\
[x_1^0]: &Q_1^1-tS_1^{-1}Q_1^2-tS_1^{-1}Q_1^0+ t \ol{x_2}Q_{1,2}^{2,0}+t\ol{x_2}Q_{1,2}^{0,0} = 0\\
                &-t\ol{x_1}S_1^{-1}Q_1^2+t\ol{x_1x_2}Q_{1,2}^{2,0} = t\ol{x_1}S_1^{-1}Q_1^0-t\ol{x_1x_2}Q_{1,2}^{0,0}-\ol{x_1}Q_1^1
\end{align}
which we'll use soon. We first start with
\begin{align}
KQ &= K_1Q_1^0 + K_2 Q_2^0-K_{1,2}Q_{1,2}^{0,0} + \ol{a_1a_2}\\ 
      &= (c_1+\wh{K}_1)Q_1^0 + K_2 Q_2^0-(c_1c_2+\wh{K}_{1,2})Q_{1,2}^{0,0} + \ol{a_1a_2}\\
K\wh{Q} &=\wh{K}_1\wh{Q}_1^0 + K_2 \wh{Q}_2^0-\wh{K}_{1,2}\wh{Q}_{1,2}^{0,0} + \ol{a_2}
\end{align}
Then we have
\begin{align}
KQ-K\wh{Q} &= \wh{K}_1(Q_1^0-\wh{Q}_1^0)+\wh{K}_2(Q_2^0-\wh{Q}_2^0)-\wh{K}_{1,2}(Q_{1,2}^{0,0}-\wh{Q}_{1,2}^{0,0})+(c_1Q_1^0-c_1c_2Q_{1,2}^{0,0}+\ol{a_1a_2}-\ol{a_2})\\
                     &=\wh{K}_1(Q_1^0-\wh{Q}_1^0)+\wh{K}_2(Q_2^0-\wh{Q}_2^0)-\wh{K}_{1,2}(Q_{1,2}^{0,0}-\wh{Q}_{1,2}^{0,0}) + tb_1S_1^{-1}Q_1^1
\end{align}
Also
\begin{align}
 \wh{K}_1(Q_1^0-\wh{Q}_1^0) -\wh{K}_{1,2}(Q_{1,2}^{0,0}-\wh{Q}_{1,2}^{0,0})=&\wh{K}_1b_1Q_1^2-\wh{K}_{1,2}b_1Q_{1,2}^{2,0}\\
=&b_1(\wh{K}_1Q_1^2-\wh{K}_{1,2}Q_{1,2}^{0,0}\\
=&b_1(-t\ol{x_1}S_1^{-1}Q_1^2+t\ol{x_1x_2}Q_{1,2}^{2,0}) \\
=&b_1(t\ol{x_1}S_1^{-1}Q_1^0-t\ol{x_1x_2}Q_{1,2}^{0,0}-\ol{x_1}Q_1^1)
\end{align}
Next
\begin{align}
K_2(Q_2^0-\wh{Q}_2^0) &=(c_2 -t\ol{x_1}^2S_2^{-1})\ol{x_2}b_1(Q_2^0-Q_{1,2}^{0,0})\\ 
                                                     &=-t\ol{x_1}^2\ol{x_2}b_1S_2^{-1}(Q_2^0-Q_{1,2}^{0,0})+c_2b_1\ol{x_1}^2(Q_2^0-Q_{1,2}^{0,0})
\end{align}
We can finally combine these expressions into
\begin{align}
x_1^2(KQ-K\wh{Q}) = b_1x_1^2(tS_1^{-1}Q_1^1+t\ol{x_1}S_1^{-1}Q_1^0-t\ol{x_1x_2}Q_{1,2}^{0,0}-\ol{x_1}Q_1^1-t\ol{x_1}^2\ol{x_2}S_2^{-1}(Q_2^0-Q_{1,2}^{0,0})+c_2\ol{x_1}^2(Q_2^0-Q_{1,2}^{0,0}))
\end{align}
We cancel out the $b_1$ and we have the equation
\begin{align}
KQ-K(x_1Q_1^1+Q_1^0)) &= tx_1S_1^{-1}(x_1Q_1^1+Q_1^0) - (tx_1\ol{x_2}Q_{1,2}^{0,0}+x_1Q_1^1+t\ol{x_2}S_2^{-1}(Q_2^0-Q_{1,2}^{0,0}))-c_2(Q_2-Q_{1,2}^{0,0})
\end{align}
Note that
\begin{align}
  K+tx_1S_1^{-1} = 1- t\ol{x_1}S_1^{-1}
\end{align}
Once we expand the $KQ$ term on the LHS we get
\begin{align}
&c_1Q_1^0-t\ol{x_1}S_1^{-1}Q_1^0+c_2Q_2^0-t\ol{x_2}S_1^{-1}Q_2^0+c_1c_2Q_{1,2}^{0,0}+t\ol{x_1x_2}Q_{1,2}^{0,0}+\ol{a_1a_2}\\
=&(1-t\ol{x_1}S_1^{-1})(x_1Q_1^1+Q_1^0)-x_1Q_1^1+t\ol{x_1x_2}Q_{1,2}^{0,0}-t\ol{x_2}S_2^{-1}Q_2^0+c_2(Q_2^0-Q_{1,2}^{0,0})
\end{align}
and after careful book-keeping all the terms cancel out one by one. This completes the long-winded proof of the identity.
%=======================================
\section{SAWs}
%=======================================

We can handle discrete self-avoiding lattice walks with this theory. If we pick a marking operator $\Gamma$ that marks every location on the plane, this will allow us to track the lattice walks' history. All we need to do in the end is remove walks that have stepped at some location more than once. However, we will investigate this problem in a bit more detail as it serves as an example of the theory. To use the generating function iteration properly we need a sort of limiting argument. For the first $n$ steps we "only" require an operator covering the square region of $(2n+1)^2$ lattice points. The overall operator is defined through a limiting process of applying the full operator (or truncated version) on a finite Laurent polynomial. Let
\begin{align}
\Gamma_{i,j} = 1+b_{i,j}\ext{x_1^ix_2^j}
\end{align}
where as usual $b_{i,j} = a_{i,j}-1$. This operator takes a generating function and marks the monomial $x_1^ix_2^j$. We will be using a lot of these together, so we let
\begin{align}
\Gamma_n = \prod_{-n\leq i,j\leq n}\Gamma_{i,j}
\end{align}
To find a functional equation for SAWs on the plane we perform an iteration. Let $Q$ be the generating function for walks with history, that is, they have a mark for each coordinate they pass through, and let $Q^{\leq n}$ be the truncated version of this generating function. It is given by
\begin{align}
Q^{\leq n} = \sum_{k=0}^{n} Q_k t^k
\end{align}
where the initial condition is $Q_0 = a_{0,0}$. Starting from $Q_n$ we can find $Q_{n+1}$ by iterating with the step set $S$, marking with $\Gamma_{n+1}$ as follows
\begin{align}
	Q_{n+1} = \Gamma_{n+1}SQ_n
\end{align}
Using a computer we can easily iterate this to add more and more steps. We can also obtain $Q^{\leq n+1}$ from $Q^{\leq n}$ since
\begin{align}\begin{split}
	t\Gamma_{n+1}SQ^{\leq n} & = \sum_{k=0}^n \Gamma_{n+1}SQ_kt^{k+1}= \sum_{k=0}^n \Gamma_{k+1}SQ_kt^{k+1} = \sum_{k=0}^n Q_{k+1}t^{k+1} = Q^{\leq n+1}-Q_0
\end{split}\end{align}
or in other words
\begin{align}
	 Q^{\leq n+1} = a_{0,0} + t\Gamma_{n+1}SQ^{\leq n}
\end{align}
In the interest of finding a functional equation we would like to take the limiit of this as $n \to \infty$. For this we define $\Gamma$ to be the operator marking all monomials $x_1^ix_2^j$, acting in the following way
\begin{align}
	\Gamma x^i y^j = \Gamma_{\max(|i|,|j|)} x^i y^j
\end{align}
and extended by linearity. The limiting process is
\begin{align}
	Q = a_{0,0} + t \Gamma SQ
\end{align}
We finally have the functional equation for lattice walks with full history
\begin{align}
	(1 -t\Gamma S)Q = a_{0,0}.
\end{align}
As in the example of an infinte marking operator, it stands to reason that $\Gamma$ is of infinite type as well. From the equation we can retrieve the generating function $W$ for the SAWs by removing all walks with power of mark $a_{i,j}$ greater than 1. For this we need an operator that removes terms with more than one copy of a given label. Let
\begin{align}
	R_{i,j} &= \ext{a_{i,j}^0} + \ext{a_{i,j}}\\
	R_n &= \prod_{-n\leq i,j \leq n} R_{i,j}\\
	R &= \lim_{n\to \infty} R_n
\end{align}
where the limit is understood as in the example of $\Gamma$ above. The solution for $Q$ can also be restated as
\begin{align}
	Q = \sum_{k=0}^{\infty} t^k (\Gamma S)^{k}(a_{0,0}) = \sum_{k=0}^{\infty} t^k (\Gamma_k S)^{k}(a_{0,0}) 
\end{align}
We can either redo the analysis, applying $R$ at every iteration, or we can apply $R$ to $Q$ and find
\begin{align}
	W = RQ = \sum_{k=0}^{\infty} t^k R(\Gamma S)^{k}(a_{0,0}) = \sum_{k=0}^{\infty} t^k (R\Gamma S)^{k}(a_{0,0}) 
\end{align}
We can write this as a functional equation
\begin{align}
	\mc L W = (1-t\Sigma)W = (1-tR\Gamma S)W = a_{0,0}.
\end{align}
It is not clear if the operator $R\Gamma$ is finite. Now it is enough to set all the remaining labels to 1 and we will find the full generating function of SAWs with steps from $\mc S $.

The discussion lends itself nicely to generalization where we have $d$ variables and operators $\Gamma_{i_1,\ldots,l_d}$, $R_{i_1,\ldots,i_d}$, and so on, defined analogously. The eigenpair discussion works out in the same way with eigenfunctions for $\lambda = 1$
\begin{align}
	u = \frac{a_{i_1,\ldots,i_d}^2 a_{j_1,\ldots,j_d}^m}{a_{k_1,\ldots,k_d}S}x_1^{r_1}\cdots x_d^{r_d} t^n
\end{align}
Also, we can constrain the SAWs on a part of $\ze^d$ with an extraction operator $E$ and mark special locations with an operator $L$. The functional equation is
\begin{align}
	(1- tLE\Sigma)W = a_{0,\ldots,0}
\end{align}
for arbitrary number of dimensions, with regional constraints and marking of special locations.

We can try finding eigenvalues and eigenfunctions of $\mc L$. Let $\mc Lu = \lambda u$. Writing it out leaves us with
\begin{align}
	t\Sigma u = (1-\lambda)u
\end{align}
So if $\lambda = 1$, $u\in \ker(\Sigma)$. That is, $R\Gamma S u = 0$. In other words, $\Gamma S u$ only has terms that involve powers greater than 1 of $a_{i,j}$ for various $i,j$. Suppose that
\begin{align}
	\Gamma S u = a_{i,j}^2 f(x_1,x_2)
\end{align}
Then
\begin{align}
	u  = \frac{a_{i,j}^2}{S}\Gamma^{\um 1}(f)
\end{align}
So we here we can let
\begin{align}
		f = a_{k,l}^m x_1^rx_2^s t^n
\end{align}
and we find that
\begin{align}
	u =  \frac{a_{i,j}^2 a_{k,l}^m}{a_{r,s}S}x_1^r x_2^s t^n
\end{align}
Thus, the kernel of $\Sigma$ is generated by these functions for all coefficients $i,j,k,l,r,s,n \in \ze$ and $m\ge 0$. This may perhaps prove useful later on. Now, if we look at eigenvalues $\lambda \neq 1$ we have to solve the equation
\begin{align}
	t \Sigma u = (1-\lambda) u
\end{align}
Because of the $t$, if $u$ is a finite order laurent series in $t$ with first nonzero coeficient $u_n$. We find that since $\Sigma$ does not interact with the $t$ monomials
\begin{align}
	\Sigma u_n = \mu u_{n-1} = 0
\end{align}
and this forces all $u_k$ coefficients to be zero. Therefore, $\Sigma$ has no other eigenvalues except 1 among the finite order power series in $t$.

%================================================================
%================================================================
%================================================================
%================================================================
%================================================================


\endinput
